\section{EPERVIER: Public Key Recovery}


\subsection{Motivation: Ethereum's Account Model}

Ethereum uses address-based accounts where:
\begin{equation}
\text{address} = \text{rightmost 20 bytes of } \text{Keccak256}(\text{pubkey})
\end{equation}

ECDSA's ecrecover precompile recovers public keys from signatures, enabling:
\begin{itemize}
    \item Transaction validation without storing public keys
    \item EIP-4337 Account Abstraction (AA) with signature aggregation
    \item EIP-7702 delegation (EOA $\rightarrow$ smart contract logic)
\end{itemize}

Standard FALCON requires explicit public key storage ($2 \times 512 = 1024$ bytes), incompatible with Ethereum's 20-byte address model.

\subsection{FALCON Recovery Mode (FalconRec)}

FALCON specification (Section 3.12) \cite{falcon2020} describes a recovery mode:

\textbf{Key Modification:}
\begin{equation}
\text{address} = H(h)
\end{equation}
where $H$ is collision-resistant hash (e.g., Keccak256).

\textbf{Signature:} $\sigma = (r, s_1, s_2)$ (includes both signature components)

\textbf{Recovery:}
\begin{algorithmic}[1]
\STATE Compute $c = \text{HashToPoint}(r \| m)$
\STATE Recover $h = s_2^{-1} \cdot (c - s_1) \mod q$
\STATE Verify $H(h) = \text{address}$
\STATE Check $\|s_1\|^2 + \|s_2\|^2 \leq B$
\end{algorithmic}

\textbf{Cost Problem:} Computing $s_2^{-1}$ in coefficient domain requires:
\begin{itemize}
    \item Extended Euclidean algorithm (expensive)
    \item Forward NTT to compute $h$
    \item Inverse NTT to convert back
    \item \textbf{Total:} 3 NTT operations (vs. 2 in standard verification)
\end{itemize}

\subsection{EPERVIER Optimization}

EPERVIER (French for ``sparrowhawk'') eliminates inverse NTT through two innovations:

\subsubsection{Hint-Based Inversion}

We provide $\text{NTT}(s_2^{-1})$ as a hint in the signature:

\textbf{Signature:} $\sigma = (r, s_1, s_2, \text{hint})$ where $\text{hint} = \text{NTT}(s_2^{-1})$

\textbf{Verification:}
\begin{algorithmic}[1]
\STATE Verify hint correctness: $\text{NTT}(s_2) \odot \text{hint} = [1, 1, \ldots, 1]$
\STATE Check norm: $\|s_1\|^2 + \|s_2\|^2 \leq B$
\STATE Compute: $h_{\text{NTT}} = \text{hint} \odot \text{NTT}(c - s_1)$
\STATE Verify: $H(h_{\text{NTT}}) = \text{address}$
\end{algorithmic}

\textbf{Cost Reduction:}
\begin{itemize}
    \item Step 1: Pointwise multiplication in NTT domain (80k gas)
    \item Step 3: One forward NTT eliminated
    \item \textbf{Result:} 2 NTT operations (same as standard FALCON)
\end{itemize}

\subsubsection{NTT-Domain Public Key}

Further optimization moves the public key itself to NTT domain:

\textbf{Address:}
\begin{equation}
\text{address} = H(\text{NTT}(h))
\end{equation}

\textbf{Recovery:}
\begin{algorithmic}[1]
\STATE $h_{\text{NTT}} \gets \text{hint} \odot \text{NTT}(c - s_1)$
\STATE Verify $H(h_{\text{NTT}}) = \text{address}$
\end{algorithmic}

\textbf{Key Insight:} By keeping computation in NTT domain, we avoid the inverse NTT entirely. The recovered value is already in the correct form for hashing.

\textbf{Final Cost:}
\begin{center}
\begin{tabular}{lrr}
\hline
\textbf{Variant} & \textbf{NTT Ops} & \textbf{Gas} \\
\hline
Standard FALCON & 2 forward, 1 inverse & 2.7M \\
FalconRec (naive) & 3 forward & 2.7M \\
EPERVIER & 2 forward & 1.9M \\
\hline
\end{tabular}
\end{center}

\subsection{Security Analysis}

\textbf{Hint Integrity:} The condition $\text{NTT}(s_2) \odot \text{hint} = [1, \ldots, 1]$ is checked in NTT domain where:
\begin{equation}
\hat{s}_2[i] \cdot \hat{h}[i] = 1 \mod q \quad \forall i \in [0, n)
\end{equation}

This is equivalent to verifying polynomial inversion in coefficient domain, since NTT is an isomorphism.

\textbf{Malleability:} FALCON signatures are naturally malleable (negation). We enforce canonical encoding by requiring all coefficients $s_{2,i} \in [0, q/2)$ (CVETH-2025-080202).

\textbf{Public Key Uniqueness:} The NTT transformation is bijective, so $H(\text{NTT}(h))$ and $H(h)$ have equivalent collision resistance. Moving to NTT domain does not weaken security.

